% Shared preamble for the English summaries in this repository.
\documentclass[11pt,a4paper]{article}
\usepackage[margin=2.2cm]{geometry}
\usepackage[T1]{fontenc}
\usepackage[utf8]{inputenc}
\usepackage{lmodern}
\usepackage{amsmath,amssymb}
\usepackage{graphicx}
\usepackage{booktabs}
\usepackage{enumitem}
\usepackage[hidelinks]{hyperref}
\setlist{itemsep=2pt,topsep=4pt}
\setlength{\parindent}{0pt}
\setlength{\parskip}{5pt}
\graphicspath{{figures/}}
\newcommand{\summaryheader}[3]{%
  \begin{center}
    {\Large\bfseries #2}\\[3pt]
    {\large #1}\\[3pt]
    {\small Arshia Goshtasbi \textbullet{} MSc Electrical Engineering (Control), Iran University of Science and Technology \textbullet{} #3}
  \end{center}
  \vspace{2pt}\hrule\vspace{8pt}
  \noindent\textit{English summary of my original Persian report, which is kept as a PDF in this folder.}\par}

\begin{document}
\summaryheader{Multivariable Control, Homework 6}{Controllability, Observability and Balanced Realisations}{Autumn 1403 (2024)}

\begin{itemize}
  \item \textbf{Near pole--zero cancellation.} I compared the step responses of $G(s)=1.01/(s+2)$ and $G(s)=(s+1.01)/((s+1)(s+2))$. The near cancellation makes the second-order system behave almost like the first-order one.
  \item \textbf{Exercises 4-1 and 4-2} (Khaki-Sedigh, \emph{Multivariable Control}). For the given system I checked with MATLAB that it is controllable, observable, functionally controllable ($\det C(sI-A)^{-1}B\neq 0$) and output controllable (rank 3). I also computed the controllability and observability indices from the matrices $\Phi_c$ and $\Phi_o$: both sets are $\{2,1,2\}$, summing to 5.
  \item \textbf{Exercises 4-3 and 4-4} from the same book.
  \item \textbf{MATLAB's Lyapunov solver.} I checked which papers the \texttt{lyap}/\texttt{dlyap} solvers are based on: Bartels--Stewart (1972), Barraud (1977), Hammarling (1982), Penzl (1998), and Golub--Nash--Van Loan (1979).
  \item \textbf{Hankel singular values.} I derived the relations between the gramians and the Hankel singular values used in balanced truncation (slide 15, session 10).
\end{itemize}
\end{document}
